\section{Introducción}
\label{sec:introduccion}

El propósito fundamental de este proyecto es el desarrollo de una solución inteligente basada en Procesamiento de Lenguaje Natural (PLN) para automatizar uno de los procesos más tediosos en la gestión financiera: la clasificación de movimientos bancarios. En un entorno digital donde el volumen de transacciones personales crece exponencialmente, la capacidad de organizar automáticamente estos datos deja de ser un lujo para convertirse en una necesidad técnica. El sistema diseñado actúa como un puente entre la información desestructurada que emiten las entidades bancarias y la estructura organizada que requiere un usuario para comprender sus hábitos de gasto.

El núcleo de la investigación se centra en la capacidad de interpretación semántica. A diferencia de los sistemas tradicionales basados en reglas rígidas o búsqueda de palabras clave —que suelen fallar ante variaciones ortográficas o sinónimos—, este sistema se enfoca en comprender la intención detrás de cada descripción. Para lograr esta robustez, el proceso de entrenamiento utiliza un corpus de datos diseñado para cubrir un amplio espectro de posibilidades lingüísticas. Se busca que el modelo no solo reconozca términos aislados, sino que sea capaz de gestionar la variabilidad con la que una misma operación financiera puede ser registrada, garantizando así una categorización coherente y precisa en escenarios de uso real.

Finalmente, el proyecto destaca por su capacidad de mejora personalizada. El sistema no es una herramienta estática, sino que permite integrar las correcciones y las nuevas descripciones que el usuario genere en su día a día. Al reentrenar el modelo con estos datos específicos, el clasificador deja de ser una solución generalista para adaptarse al vocabulario y a los hábitos de gasto particulares de cada persona. De esta forma, el sistema garantiza una precisión creciente, optimizando su rendimiento a medida que se acumula un historial de transacciones real y verificado por el propio usuario.